% !TEX root = tarasov-qasp-zh.tex


\chapter[你会作平面镜和透镜所成的像吗？]{你会作平面镜和透镜所成的像吗？\label{ch-32}}

\begin{dialogue}

\tchr 我们常常发现，考生不会作各种光学系统——例如透镜、平面镜（plane mirror）和球面镜（spherical mirror）——所成的像。下面来看几个典型的例子。\hlred{作出人在\figr{fig-132} 所示平面镜中所成的像。}

\begin{figure}%
\includegraphics[width=.7\textwidth]{figs/ch-32/fig-132a}
\caption[作出人在平面镜中所成的像。]{作出人在平面镜中所成的像。\label{fig-132}}
\end{figure}

\stda 我觉得这种情形下镜中不会成像，因为镜子在人头顶上方挂得太高了。

\tchr 你错了。镜中是会有像的。作图见图\figr{fig-132}~\drkgry{(b)}。十分明显，要作出这个像，只需把表示镜面的直线延长，再画出人形关于这条直线（镜面）的对称图形。

\stda 是的，这我明白了。可是，那个人能看见自己的像吗？

\tchr 这就是另一个问题了。事实上，那个人是看不见自己的像的：镜子在他头顶上方太高，倾角也不合适。在给定的这面镜子里，只有位于光线 $AA_{1}$ 与 $BB_{1}$ 所夹角度内的观察者才能看到人的像。这里正好回忆一下：观察者的眼睛从被观察的物体接到的是一束发散光线。眼睛将在这些光线或其反向延长线的交点处看到物体的像（见\figr{fig-129} 与\figr{fig-132}~\drkgry{(b)}）。

\hlred{现在来考虑由两块互相垂直放置的平面镜组成的系统所成的像的作图（\figr{fig-133}~\drkgry{(a)}）。}
\begin{figure}%
\centering
\includegraphics[width=.7\linewidth,angle=-3]{figs/ch-32/fig-133a}
\caption[作图求两块互成直角的平面镜所成的像。]{作图求两块互成直角的平面镜所成的像。\label{fig-133}}
\end{figure}
\stda 只要作出物体在两个镜面里的反射像就行了。这样就得到两个像，如\figr{fig-133}~\drkgry{(b)} 所示。

\tchr 你把第三个像丢了。注意，从物体发出、位于直角 $AOB$ 之内（\figr{fig-133}~\drkgry{(c)}）的那些光线要经受两次反射：先从一面镜反射，再从另一面镜反射。\figr{fig-133}~\drkgry{(c)} 中画出了这两条光线的路径。这些光线反向延长线的交点就定出物体的第三个像。
\pnote{两块平面镜成 $\theta$ 角组成“角镜”。当 $360^{\circ}/\theta$ 为偶数时（$\theta=90^{\circ}$ 即是），所成的像共 $360^{\circ}/\theta-1$ 个：$\theta=90^{\circ}$ 给出 3 个像，其中第三个由先后经两镜各反射一次的光线形成，位置恰与物关于两镜的交点中心对称。角镜对光线的这种“折返”作用后来被用于角反射器：三面互成直角的镜面使入射光沿原方向返回，自行车尾灯与卫星激光测距都利用了这一点。}

下面我们来考虑几个与会聚透镜（converging lens）有关的例子。试就\figr{fig-134}~\drkgry{(a)} 所示的情形，作出这种透镜所成的像。

\stda 这很简单。我的作图如\figr{fig-134}~\drkgry{(b)} 所示。

\begin{figure*}%
\includegraphics[width=0.8\textwidth]{figs/ch-32/fig-134a}
\caption[会聚透镜的几个例子。]{与会聚透镜有关的几个例子。\label{fig-134}}
\end{figure*}

\tchr 很好。现在假定透镜的一半被一块不透明屏挡住，如\figr{fig-134}~\drkgry{(c)} 所示。像会变成什么样？
\pnote{平面镜成像与凸透镜成像规律在现行课程中属初中内容：人教版八年级上册第四章《光现象》、第五章《透镜及其应用》；原书所说的会聚透镜、发散透镜即初中教材的凸透镜、凹透镜。高中选择性必修第一册第四章《光》讲折射、全反射与干涉、衍射、偏振，不再系统训练透镜成像作图。本章后文的焦平面作图法在中学教材中不出现，是物理竞赛里的常用工具。}

\stda 这种情形下像就没有了。

\tchr 你错了。你忘了：箭头上任意一点（比如箭头尖端）的像，是由无穷多条光线相交得到的（\figr{fig-134}~\drkgry{(d)}）。我们通常只限于两条光线，是因为要用作图法定出像的位置和大小，两条光线的路径就足够了。在本题的情形里，屏遮住射向透镜的一部分光线；然而其余部分的光线仍将穿过透镜并形成物体的像（\figr{fig-134}~\drkgry{(e)}）。由于参与成像的光线减少了，像将不如先前那么亮。

\stdb 由您的解释可以得出：用不透明屏遮住透镜的一部分时，改变的只是像的亮度，其他什么也不变。可是，凡是同摄影打过交道的人都知道：收缩相机的光圈（aperture）以减小其孔径，也就是减小透镜的有效面积时，除了像的亮度降低之外，还观察到另一种效果——像变得更清晰、更分明了。为什么会这样？

\tchr 这个问题提得很中肯。它使我可以强调下面这一点：我们前面全部作图所依据的假定，是可以忽略光学系统（本例中就是透镜）的缺陷。的确，“缺陷”一词在这里并不十分恰当，因为它涉及的并不是透镜偶然的缺点，而是透镜的基本性质。大家知道，若有两条都平行于主光轴（principal optical axis）、但离轴距离不同的光线穿过透镜，那么严格说来，它们经透镜折射后将与主光轴交于不同的点（\figr{fig-135}~\drkgry{(a)}）。这就是说，透镜的焦点（focus）——所有平行于主光轴的光线的交点——将是“弥散的”，不可能形成轮廓清晰的像。各条光线离主光轴的距离差别越大，像就越模糊。收缩光圈后，透镜将通过一束狭窄的光线，这在一定程度上改善了清晰度（\figr{fig-135}~\drkgry{(b)}）。
\pnote{原书所说的透镜“缺陷”，现代术语叫像差（aberration），此处描述的是其中最基本的一种——球面像差（spherical aberration，球差）：平行于主光轴的光线随离轴高度不同而交轴于不同点，\figr{fig-135} 中的 $\Delta F$ 正标出了这一差值。“平行光线交于一点”只在近轴条件（只考虑靠近光轴的细光束）下成立——这是一切成像作图的前提；中学教材默认这一近似而不提像差，到大学光学课程才系统处理。另外，光圈也不能一味收小：孔径太小时衍射的影响上升，像反而重新变模糊，故存在最佳光圈。}
\begin{figure}%
\centering
\includegraphics[width=\textwidth]{figs/ch-32/fig-135a}
\caption[透镜的成像作图。]{透镜的成像作图。\label{fig-135}}
\end{figure}

\stdb 这么说，借助光圈，我们是以牺牲亮度为代价，使像的轮廓变得更清晰的。

\tchr 正是。不过请记住：在作透镜成像图时，考生完全有理由认为平行光线总会交于一点。如果平行光束沿主光轴方向射来，这个点就在主光轴上；如果平行光束与主光轴成某一角度射来，这个点就在焦平面（focal plane）上。然而同样重要的是，考生要明白：这种处理只是近似的，更精确的方法须对光学系统的缺陷加以修正。

\stda 什么是透镜的焦平面？

\tchr 它是通过透镜主焦点并垂直于主光轴的平面。那么，在\figr{fig-134} 的例子中，平面镜所成的像与会聚透镜所成的像有什么差别？

\stda 第一种情形（平面镜）的像是虚像（virtual image），第二种情形的像是实像（real image）。

\tchr 正确。请更详细地解释虚像与实像的差别。

\stdb 虚像不是由光线本身、而是由光线的反向延长线相交形成的。难怪虚像会出现在墙后之类光线根本进不去的地方。

\tchr 完全正确。还要注意，虚像只能从一些确定的位置观察到；而对于实像，你可以在成像处放一个屏，从任意位置观察它。来看\figr{fig-136}~\drkgry{(a)} 所示的例子。

\hlred{已知另一条光线（\figr{fig-136}~\drkgry{(a)} 中的 $BB_{1}B_{2}$）经过会聚透镜的路径，试用作图法确定光线 $AA_{1}$ 通过该透镜后的方向。}

\begin{figure}[h]%
\centering
\includegraphics[width=0.5\textwidth]{figs/ch-32/fig-136a}
\caption[已知一条光线的路径，作图求另一条光线的方向。]{已知光线 $BB_{1}B_{2}$ 经会聚透镜的路径，用作图法求光线 $AA_{1}$ 通过透镜后的方向：(b)、(c) 为辅助箭头法，(d) 为焦平面法。\label{fig-136}}
\end{figure}
\pnote{底本\figr{fig-136} 的图题误沿了\figr{fig-133} 的文字（``Construct the images in the system of two plane mirror''），与本图内容（会聚透镜的光路作图）不符，显系录入之误；中文版按图面与正文拟写图题，待考原书。下文教师答语中，底本又有整句误重（``we can determine the direction of the ray $AA_{1}$ through the lens, we can determine the direction of ray $AA_{1}$ after passing through the lens''），按上下文只取一次，中文照此译出。本章另有多处拼写与标点级小误（shown、from、图题名词单复数、图号引用符错位等），径改不注。}

\stda 可是，我们并不知道透镜的焦距（focal length）呀。

\tchr 可我们知道另一条光线在透镜之前和之后的路径。

\stda 这种作图我们在学校里没有学过。

\stdb 我认为应当先求出透镜的焦距。为此，可以在透镜左边任取一处画一个竖直的箭头，使箭头的尖端与光线 $BB_{1}$ 相接触。把箭头尖端记作字母 $C$（\figr{fig-136}~\drkgry{(b)}）。然后从 $C$ 点引一条经过透镜中心的光线。这条光线将不发生折射地穿过透镜，并在某点 $E$ 与光线 $B_{1}B_{2}$ 相交。点 $E$ 显然就是箭头尖端的像。剩下的事，是从箭头尖端 $C$ 引第三条光线，使它平行于透镜的主光轴。这最后一条光线折射后将经过箭头尖端的像，即经过点 $E$。这第三条光线与主光轴的交点，就是所求的透镜焦点。这个作图见\figr{fig-136}~\drkgry{(b)}。

焦距既已知道，现在就可以作出光线 $AA_{1}$ 被透镜折射后的路径了。做法是再画一个竖直箭头，使其尖端位于光线 $AA_{1}$ 上（\figr{fig-136}~\drkgry{(c)}）。利用已确定的焦距，可以作出第二个箭头的像。所求的光线将经过点 $A_{1}$ 以及箭头之像的尖端。这个作图如\figr{fig-136}~\drkgry{(c)} 所示。

\tchr 你的论证完全正确，其依据是求出某个辅助物体（箭头）的像。注意，当要求确定位于透镜主光轴上的发光点（luminous point）的像的位置时，这个方法是很方便的：这时宜在发光点处竖立一个箭头，作出箭头的像。显然，箭头之像的尾端就是所求的发光点的像。

不过，对我们这个例子来说，这个方法太繁琐了。我来讲一个更简单的作图。为求出透镜的焦距，可以经过透镜中心作一条平行于光线 $BB_{1}$ 的光线 $EO$（\figr{fig-136}~\drkgry{(d)}）。由于这两条光线平行，它们经透镜折射后将交于透镜后方焦平面内的一点（\figr{fig-136}~\drkgry{(d)} 中用虚线示出了焦平面的截线）。然后经过透镜中心作平行于光线 $AA_{1}$ 的光线 $CO$。既然这两条平行光线过透镜后同样应当交于焦平面内，我们就能确定光线 $AA_{1}$ 通过透镜之后的方向。可以看到，这个作图要简单得多。

\stdb 是的，您的方法明显地简单一些。

\tchr 试着把这个方法用到一道类似的题目上，那里用发散透镜（diverging lens）代替了会聚透镜（\figr{fig-137}~\drkgry{(a)}）。
\begin{figure}%
\centering
\includegraphics[width=0.5\textwidth]{figs/ch-32/fig-137a}
\caption[用发散透镜作同样的作图。]{用发散透镜作同样的作图。\label{fig-137}}
\end{figure}
\stdb 我先经过透镜中心作一条平行于光线 $BB_{1}$ 的光线。与上一题不同，这里相交的将是光线的反向延长线，而不是光线本身（可以指出，对经过中心的光线来说，反向延长线与光线本身重合）。于是，包含交点的焦平面现在将位于透镜的左侧而不是右侧（见\figr{fig-137}~\drkgry{(b)} 中的虚线）。

\tchr （插话）注意，发散透镜所成的像总是虚像。

\stdb （继续）接下来我再经过透镜中心作一条平行于光线 $AA_{1}$ 的光线。从这两条光线的反向延长线交于焦平面内这一条件出发，我就能画出所求的光线。

\tchr 好。现在告诉我：一个物体，它的一部分位于会聚透镜的焦点以内，另一部分位于焦点以外（物体具有一定的宽度），它的像在哪里？

\stdb 我将作出物体上若干个离透镜距离不同的点的像。位于焦点以外的点将给出实像（它位于透镜右侧）；位于焦点以内的点将给出虚像（它位于透镜左侧）。当所取的点趋近焦点时，像将向无穷远处移去（向透镜左侧或右侧）。

\tchr 好极了。这样，在我们这个情形里，物体的像由两段组成（分别在透镜左侧和右侧）。每一段都从离透镜某个有限距离处开始，一直延伸到无穷远。正如你们所看到的，对“一个物体能不能同时具有实像和虚像？”这个问题，应当作出肯定的回答。

我看得出来，你们已经掌握了透镜成像的作图步骤。因此可以转入更复杂的题目——由两个透镜组成的系统所成像的作图。考虑下面的问题：

\hlred{有两个共主光轴而焦距不同的会聚透镜。试作出竖直箭头经这样的光学系统所成的像\drkgry{（\figr{fig-138}~(a)）}。图中一个透镜的焦点用叉号 $\times$ 标出，另一个透镜的焦点用实心黑圆点 $\bullet$ 标出。}

\begin{figure}%
\centering
\includegraphics[width=0.5\textwidth]{figs/ch-32/fig-138a}
\caption[作图求光学系统所成的像。]{作图求两个透镜组成的光学系统所成的像。\label{fig-138}}
\end{figure}
\stdb 要作出箭头经两个透镜所成的像，必须先作出第一个透镜所成的像，这时可以不管第二个透镜；然后把这个像当作一个物体，再不管第一个透镜，作出它经第二个透镜所成的像。

\tchr 你这里犯了一个非常典型的错误。这样的回答我已听到过许多次。它是完全错误的。

我们来考察从箭头尖端发出的两条光线，跟踪它们通过给定透镜系统的路径（\figr{fig-138}~\drkgry{(b)}）。这两条光线过第一个透镜后的路径是容易作出的。为了找出它们过第二个透镜后的路径，我们要画出与它们平行、并经过第二个透镜中心的辅助光线。这时我们利用了前面几题中讨论过的原理（平行的光线经透镜后应交于焦平面内）。所求的箭头尖端的像，就在两条原始光线离开第二个透镜后的交点上。这个作图详细地画在\figr{fig-138}~\drkgry{(b)} 中。现在来看一看，如果接受你的建议，会得到什么结果。作图画在\figr{fig-138}~\drkgry{(c)} 中：实线表示第一个透镜成像的作图，虚线表示随后第二个透镜成像的作图。你们可以看到，那样得到的结果会完全不同（而且完全不对！）。

\stdb 可是我确信，我们有一次就是按我指出的那样作过像的。

\tchr 你可能那样做过。事实是，在某些情形里，你的作图方法可能恰好成立，因为它给出的结果与我的方法给出的结果相同。这可以在前一个例子上加以演示：把箭头移近第一个透镜，即移到焦点与透镜之间。\figr{fig-139}~\drkgry{(a)} 画出的是按我的方法所作的图，\figr{fig-139}~\drkgry{(b)} 是按你的方法所作的图。你们看，在这种情形里两者的结果一致。
\pnote{原书判为“完全错误”的，是不加辨析地把中间像当作实物、套用特殊光线的作法。严格说来，“逐次成像”——把前一透镜的像当作后一透镜的物——在近轴光学中总是成立的，关键是区分实物与虚物：本例中间像落在第二个透镜之后，射到第二镜的是向该点会聚的光束，此“物”为虚物（透镜公式中其物距取负），仍把它当实物、从该点“引出”发散的特殊光线作图，必然得到\figr{fig-138}~\drkgry{(c)} 那样的错误结果；中间像位于第二镜入射光一侧（两镜间的实像，或\figr{fig-139} 情形的物侧虚像）时则是实物，逐次成像与逐点追迹结果一致。}

\begin{figure}%
\centering
\includegraphics[width=0.5\textwidth]{figs/ch-32/fig-139a}
\caption[两种作图方法结果一致的情形。]{把箭头移到第一个透镜的焦点与透镜之间后，(a)、(b) 两种作图方法的比较——两者结果一致。\label{fig-139}}
\end{figure}
\pnote{底本\figr{fig-139} 的图题同样误沿了\figr{fig-133} 的文字（``two plane mirror''），与本图内容（两种作图方法的比较）不符，显系录入之误；中文版按图面与正文拟写图题，待考原书。}

\stdb 可是，怎样才能事先确定，在哪些情形下可以用我的成像作图方法呢？

\tchr 对两个透镜的情形，给出你的方法的适用条件并不困难；透镜个数更多时，这些条件会变得复杂得多。根本用不着去讨论它们。用我的方法，你就不会陷入困境。不过我想再问一个问题：双凹透镜（double-concave lens）能不能是会聚的？

\stdb 在通常条件下，双凹透镜是发散的。可是，若把它放入折射率高于透镜材料折射率的介质中，它就会成为会聚透镜。在同样条件下，双凸透镜（double-convex lens）则会成为发散透镜。
\pnote{透镜会聚还是发散，取决于透镜与介质折射率的相对大小：透镜折射率较大时，双凸透镜会聚、双凹透镜发散，较小时相反；初中教材的说法默认透镜处在空气中。水中气泡形如“双凸透镜”却对光发散——竞赛里常见的概念题。}


\end{dialogue}


\cleardoublepage

% NEWTERMS: plane mirror | 平面镜 | 首现括注
% NEWTERMS: spherical mirror | 球面镜 | 首现括注；本章仅开篇列举，正文未展开
% NEWTERMS: converging lens | 会聚透镜 | 首现括注；即人教版初中的凸透镜（见评注）
% NEWTERMS: diverging lens | 发散透镜 | 首现括注；即人教版初中的凹透镜（见评注）
% NEWTERMS: virtual image | 虚像 | 首现括注
% NEWTERMS: real image | 实像 | 首现括注
% NEWTERMS: focus / focal point | 焦点 | 首现括注；principal focus（主焦点）仅出现一次，未括注
% NEWTERMS: focal length | 焦距 | 首现括注
% NEWTERMS: focal plane | 焦平面 | 首现括注；中学教材不出现（见评注）
% NEWTERMS: principal optical axis | 主光轴 | 首现括注
% NEWTERMS: luminous point | 发光点 | 首现括注
% NEWTERMS: index of refraction | 折射率 | 承 ch-31 定名（原书符号 $n$，人教版定义 $n=c/v$），ch-31 已先现，本章不重复括注
% NEWTERMS: double-concave lens | 双凹透镜 | 首现括注
% NEWTERMS: double-convex lens | 双凸透镜 | 首现括注
% NEWTERMS: aperture / irising / diaphragm | 光圈；收缩光圈 | 首现括注（aperture）
% NEWTERMS: optical system | 光学系统 | 未括注
% NEWTERMS: opaque screen | 不透明屏 | 未括注
% NEWTERMS: extension(s) of rays | （光线的）反向延长线 | 光学惯用表述，未括注
% NEWTERMS: aberration / spherical aberration | 像差；球面像差（球差） | 仅评注层引入
% NEWTERMS: virtual object | 虚物 | 仅评注层引入
% NEWTERMS: angle mirror | 角镜 | 仅评注层引入
% NEWTERMS: corner reflector | 角反射器 | 仅评注层引入
